% =====================================================================
\section{Conclusion and Future Work}\label{sec:conclusion}
% =====================================================================

We introduced \textbf{Walk Attention}, an architecture that mitigates
over-squashing in graph neural networks by attending over a multi-hop support
supplied by the path algebra of the underlying quiver. The graded path algebra
gives, through Proposition~\ref{prop:count}, an exact account of which node pairs
are connected by length-$k$ walks and with what multiplicity; we use this to
define a sparse attention whose support is algebraically determined and whose
weights are learned. On a tunable bottleneck benchmark, Walk Attention solved a
long-range retrieval task exactly ($1.000\pm0.000$ over five seeds) where a
one-hop attention network collapsed to chance and a fixed-weight multi-hop
aggregator only partially succeeded; on the real Peptides-func benchmark it was
competitive with strong baselines while remaining sparse and stable. Its
advantage is regime-dependent---decisive when over-squashing is severe and
comparable to one-hop attention when it is mild. A complementary negative result
showed that using the path-algebra \emph{quotient} to compress equivalent walks
is counterproductive when path multiplicity carries signal; the role of the
algebra is to determine the attention support, not to discard information.

The construction is deliberately elementary and self-contained, so that it can
serve as an entry point connecting quiver representation theory to graph
representation learning. Walk Attention is the learned, representational reading
of a common algebraic substrate (Section~\ref{sec:ground})---the graded path
algebra and its impact-degree neighbourhood system---whose dynamical reading is
the automata-of-impact-over-quivers framework \cite{zainea2026aiq}; relating the
two is itself a direction. Several others follow naturally.

\begin{itemize}
  \item \textbf{From learned weights to dynamics.} The dynamical sibling
  \cite{zainea2026aiq} fixes the node-update rule $\varphi$ \emph{a priori},
  weighting influence by impact degree; Walk Attention instead \emph{learns}
  those weights. This suggests a two-way bridge: the attention weights learned on
  a task can be injected into an AIQ as a data-driven evolution rule, letting one
  read off the \emph{dynamical} effect of each weight---how strongly a given
  source, at a given impact degree, drives the state of a node over time---and,
  conversely, letting AIQ stability and reachability analyses interpret and
  constrain what Walk Attention has learned. Quantifying this correspondence
  between learned attention and impact-weighted dynamics is, to us, the most
  promising direction.
  \item \textbf{Benchmarks with severe over-squashing.} Peptides-func is a mild
  case where the multi-hop support is not decisive. Identifying real long-range
  tasks whose graphs do contain acute bottlenecks---where the synthetic advantage
  is expected to carry over---and comparing against rewiring
  \cite{topping2022understanding} and graph-Transformer baselines, is the natural
  next step.
  \item \textbf{Quotient-shaped supports.} The quotient $\kQ/I$, rather than
  compressing the signal, can \emph{partition} walks into equivalence classes;
  routing distinct classes to distinct attention heads would let the algebra
  structure the heads themselves. Identifying which relations $I$ are useful---and
  learning them from data, much as the dynamical sibling fixes them by a rule---is
  an appealing problem.
  \item \textbf{Scaling the support.} For dense graphs the range-$k$ support can
  itself grow large; truncating, sampling, or learning a sparser admissible
  support are natural ways to control cost while preserving long range.
\end{itemize}

\subsubsection*{Reproducibility.}
All graph generators, the path-algebra walk operators, the Walk Attention layer,
the baselines, and the experiment-tracking logs are released
\cite{repo_walkattention}. The path-algebra computations build on the open-source
\texttt{aiq-quivers} library \cite{repo_aiqquivers}; the present work is, to our
knowledge, the first published application of that algebraic machinery to graph
neural networks, and we have therefore kept the algebraic construction explicit
throughout.
